\begin{enumerate}
	\item \begin{itemize}
		\item On choisit 5;

		\item d'une part, le carré de 5 est :\quad $5^2 = 25$;

		puis, en multipliant ceci par 3 :\quad $3 \times 25 = 75$;

		\item d'autre part, en multipliant 5 par 7 :\quad $5\times 7 = 35$;

		\item on additionne les deux nombres :\quad $75+35=110$;

		\item Finalement, on soustrait 6 :\quad $110-6=104$.
	\end{itemize}
	Le \textbf{programme A}, appliqué à 5, donne donc 104.

	\item On veut dans \textsf{B2} le résultat du \textbf{programme A}, appliqué au nombre qui est dans la cellule \textsf{A2} :
	\begin{itemize}
		\item On choisit \textsf{A2};

		\item d'une part, le carré de \textsf{A2} est \textsf{A2} multiplié par lui même:\quad \textsf{A2*A2};

		puis, en multipliant ceci par 3 :\quad \textsf{3*A2*A2};

		\item d'autre part, en multipliant \textsf{A2} par 7 :\quad $7*\textsf{A2}$;

		\item on additionne les deux nombres :\quad \textsf{3 * A2 * A2 + 7 * A2};

		\item Finalement, on soustrait 6 :\quad \textsf{3 * A2 * A2 + 7 * A2 - 6};
	\end{itemize}
	Le \textbf{programme A} sera appliqué par la formule :\textsf{3 * A2 * A2 + 7 * A2 - 6}.

	\item À l'aide du tableur, on constate que lorsqu'on prend $-3$ comme nombre de départ (cellule \textsf{A3}), on obtient 0 (cellule \textsf{B3}).

	\item \begin{itemize}
		\item On choisit $x$;

		\item d'une part, le carré de $x$ est :\quad $x^2$;

		puis, en multipliant ceci par 3 :\quad $3x^2$;

		\item d'autre part, en multipliant $x$ par 7 :\quad $7x$;

		\item on additionne les deux nombres :\quad $3x^2+7x$;

		\item Finalement, on soustrait 6 :\quad $3x^2+7x-6$.
	\end{itemize}

	Avec $x$ comme nombre de départ, une expression littérale du \textbf{programme A} en fonction de $x$ est:\quad $3x^2+7x-6$

	\item \begin{itemize}
		\item On choisit 5;

		\item d'une part, en multipliant 5 par 3 :\quad $3 \times 5 = 15$,

		puis, en soustrayant 2 : \quad $15 - 2 = 13$;

		\item d'autre part, en ajoutant 3 à 5 :\quad $5+3=8$;

		\item on multiplie les deux nombres :\quad $13 \times 8 = 104$;
	\end{itemize}

	En appliquant le \textbf{programme B} à 5, on obtient aussi 104.

	\item \begin{itemize}
		\item On choisit $x$;

		\item d'une part, en multipliant $x$ par 3 :\quad $3x$,

		puis, en soustrayant 2 : \quad $3x-2$;

		\item d'autre part, en ajoutant 3 à $x$ :\quad $x+3$;

		\item on multiplie les deux nombres :\quad $(3x-2)(x+3)$;
	\end{itemize}

	En appliquant le \textbf{programme B} à $x$, on obtient l'expression :\quad $(3x-2)(x+3)$.

	\item Trouvons la forme développée réduite de l'expression littérale du \textbf{programme B} :
	\begin{align*}
	(3x-2)\times (x+3)&=3x\times x+3x\times 3+(-2)\times x+(-2)\times 3\\
	&=3x^2+9x-2x-6\\
	&=3x^2+7x-6
	\end{align*}

	On retrouve la même expression que le \textbf{programme A}, donc Mathis a raison.

	\item Résolvons :
	\begin{center}
	$(3x-2)(x+3) = 0$

	Un produit est nul si et seulement si l'un de ses facteurs est nul ; donc :

	$3x-2=0$ \quad ou \quad $x+3=0$

	$3x=2$\quad ou \quad $x=-3$

	$x=\dfrac{2}{3}$\quad ou \quad $x=-3$
	\end{center}

	Donc l'équation a deux solutions, son ensemble des solutions est :\quad $\mathcal{S}=\left\{-3 ;\dfrac{2}{3} \right\}$.


	Ainsi les valeurs de $x$ pour lesquelles les \textbf{programmes A} et \textbf{B} donnent 0 sont $-3$ et $\dfrac{2}{3}$.
\end{enumerate}
